\section{Timer and CLINT}

\subsection{Why a timer matters}

A timer interrupt lets software regain control periodically. This is essential for:

\begin{itemize}
\item preemptive multitasking;
\item RTOS tick interrupts;
\item timeouts;
\item scheduling;
\item basic OS services.
\end{itemize}
Without a timer, one program can run forever unless it voluntarily yields.

\subsection{Machine timer interrupt flow}

The timer peripheral asserts a timer interrupt level. The CSR unit checks:

\begin{lstlisting}[language={}]
mstatus.MIE && mie.MTIE && timer_irq
\end{lstlisting}

If true, \texttt{irq\_pending} becomes active and the CPU takes a machine timer interrupt. The CSR unit reports cause \texttt{0x8000\_0007} for timer interrupt. External interrupts, such as UART-style interrupts in later variants, use cause \texttt{0x8000\_000B} and have priority in \texttt{csr.v}.

\subsection{CLINT-style block}

\texttt{rtl/clint.v} is the CLINT-style timer block used by RTOS-oriented SoC variants. In RISC-V systems, CLINT normally contains \texttt{mtime} and \texttt{mtimecmp}. When \texttt{mtime >= mtimecmp}, a timer interrupt becomes pending.

The important concept is that software programs a compare value in memory-mapped registers, then the hardware raises an interrupt when time reaches that value.

\bigskip
